% conclusions.tex (v11)
\section{Conclusions}\label{sec:conclusions}
We evaluated active-layer thickness models in a site-year design, in which labels, climate inputs and scores refer to the same observation site and year. Matching the climate input to the observation year reduced the spatial-block error where the observations span many years and assigns every prediction to a stated year. A Stefan model recalibrated with local labels and corrected by a gradient-boosting model of its residual reduced the Stefan error by \ndrefLeblkRoabs{} cm in the Lena Delta and \ndrefCablkRoabs{} cm in Canada, remained close to the Stefan model in a region withheld from training, and reduced its error in the Lena Delta and the external CALM regions with 5--25 local labels. Direct ML reached lower point errors within regions with all labels, but without Lena Delta labels its error there exceeded that of the Stefan model by \ndrefLeregDzabs{} cm. In the same cells and years, the model was more accurate than five published products in a region that the products did not learn and performed on par with Wei et al. and Liu et al. where these products were trained on nearby observations. Physical anchoring of the learner is therefore most useful where target-region labels are few or absent, which is the common situation in permafrost regions. The workflow provides year-labelled 1 km ALT maps with an area of applicability and prediction intervals, together with a protocol for evaluating and adapting ALT models and products in other permafrost regions.
